\section{Two Execution Contracts}
\subsection{A Common Native Task Boundary}
The native builder enumerates the original samples and supplies source RR identity, source-relative cell mapping, target-layer identity, CSR connectivity and indexed costs. Each backend returns sample cost pairs to the outer builder. The host retains group commit, interpolation and source/OPIN/minimum work. Both paths use this interface, but their search order and within-sample representative selection differ. This distinction prevents a shared interface from being mistaken for shared semantics.

Source identity is essential. Consider starts at $x=0$ and $x=5$ reaching a pin at $x=6$ with delays 1 and 2. They populate displacement buckets 6 and 1. A multi-source search retaining only the cheaper pin label loses the second contribution. We therefore preserve separate logical samples, including their coordinate contexts. Allocation reuse is allowed; reusing a previous sample's mutable labels or cached answer is not.

\begin{figure}[t]\centering
\begin{tikzpicture}[>=Stealth, every node/.style={draw,rounded corners,align=center,font=\footnotesize}]
\node[text width=7.4cm] (a) at (0,0) {Native sample identity, layer-relative cells, graph and indexed costs};
\node[text width=2.1cm] (s) at (-2.55,-1.25) {STRICT CPU\\private floods\\ordered reduction};
\node[text width=2.1cm] (g) at (0,-1.25) {ALT CPU/CUDA\\whole-pair\\bucket waves};
\node[text width=2.1cm] (f) at (2.55,-1.25) {Finite ALT RTL\\bounded state\\memory service};
\node[text width=7.4cm] (o) at (0,-2.5) {Native outer construction: group commit, interpolation, source/OPIN and minima};
\draw[->] (a)--(s);\draw[->] (a)--(g);\draw[->] (a)--(f);
\draw[->] (s)--(o);\draw[->] (g)--(o);\draw[->] (f)--(o);
\end{tikzpicture}
\caption{A common sample boundary with distinct semantic and implementation obligations. The FPGA outer interface is evaluated by replay and simulation, not board-integrated native timing.}\label{fig:contracts}
\end{figure}


\subsection{STRICT: Private Floods, Ordered Reduction}
Each worker owns a visited-cost array, expanded bitmap, priority queue and sample output. Read-only graph/cost data are shared. The workspace is reset before every sample. Groups preserve X/Y/Z channel filtering and all destination-layer slices. Sample search uses the native full-device bounding box and existing edge/cost operations.

The native smallest-entry rule replaces a populated cell only for a strictly smaller delay. Congestion belongs to the winning delay; it must not be minimized independently. For finite occupied representatives, the winning sample is the earliest minimum-delay sample in native enumeration order:
\begin{equation}
j^*(u)=\arg\min_{j\,\mathrm{occupied}}(d_{g,j}(u),j).
\end{equation}
Implementation uses the native insertion and occupancy predicates, including missing-state behavior, rather than generic tuple sorting.

\begin{proposition}
If native sample selection and within-flood visitation are unchanged, valid recorded delays are finite, and final commits/interpolation follow native order, private sample execution followed by original-order strict-delay reduction yields the native wire-cost table.
\end{proposition}
\noindent\emph{Proof.} Each private flood produces the same per-cell representative as its native flood. Serial execution concatenates the visitation sequences in sample order. Original-order reduction selects the minimum delay and retains the earliest associated congestion on ties. Occupancy agrees cell by cell; identical commits and interpolation preserve that agreement. Task completion order does not enter the argument. $\square$

The proposition does not permit arbitrary exceptional-value reassociation, changed heap tie order or altered recording contexts. It also concerns the wire map, not every derived source/OPIN object or eventual routing trajectory. These boundaries matter when comparing independently repaired executables.

Group executes samples concurrently but waits at each group boundary. Pipeline exposes all groups through a bounded stream with parallel search and original-order reduction; reusable private workspaces avoid repeated allocation. Flat is an independent globally flattened parallel-loop control, not VPR's separate flat-routing option. Distinct layer/type/segment slices permit disjoint work, but interpolation within a slice follows its native dependency/order. Compact storage and alternative commit schedules are secondary variants, not assumed prerequisites of the reported native GPU result.

\subsection{ALT: Whole-Pair FP32 Bucket Waves}
ALT\_BUCKET\_FP32\_V1 is a separate contract. Admitted costs and accumulated candidates must be positive normal FP32 values or positive zero (negative zero and subnormals are rejected); unsupported numerics or capacity failures return an error without publishing a partial map. The unchanged native admission includes $N<2^{23}$. Labels begin at $(+\infty,+\infty)$ except source $(0,0)$; missing output cells retain their designated missing state.

For labels $a,b$, strict improvement means
\begin{equation}
a\prec b\iff a_d<b_d\ \lor\ (a_d=b_d\land a_q<b_q).
\end{equation}
The entire pair is selected. For nonnegative finite floats, CUDA packs the delay bits above congestion bits in a 64-bit word, allowing an unsigned atomic minimum to select one actual pair. It never combines a delay from one path with congestion from another. FP32 additions use round-to-nearest with contraction disabled in the matched implementations.

Scheduling uses $B(d)=\mathrm{bits}_{32}(d)\gg s$. We fix $s=20$ for the native experiments. Shift coarsens scheduling buckets, not stored delays: labels and returned costs remain FP32. At shift zero, the schedule key is the exact nonnegative FP32 delay bit key; this fact alone does not make ALT equivalent to native tie semantics.

A wave selects the smallest pending bucket, filters stale notifications and takes the current labels as its snapshot. All eligible outgoing edges generate proposals from that snapshot. Per-destination proposals begin from the existing label and reduce whole pairs; labels are committed only after expansion completes. Strictly improved destinations are published to their new bucket, including reinsertion into the current bucket. Pending deduplication avoids losing same-bucket notifications; old notifications are filtered on extraction. Closure occurs only when no pending work remains in the selected bucket. Target cells are recorded after closure, using whole-pair selection, and completion is tested then.

\begin{algorithm}[t]\caption{ALT wave contract (one logical sample)}\label{alg:alt}
\begin{algorithmic}[1]
\State Reset labels, pending state and outputs; publish source
\While{pending work exists and target cells incomplete}
\State Select minimum pending bucket $b$; filter stale members
\State Read wave-start labels; generate eligible proposals
\State Reduce each destination by whole-pair $\prec$
\State Commit strict improvements; publish or deduplicate
\If{bucket $b$ has no remaining work}
\State Record closed target cells; test target completion
\EndIf
\EndWhile
\State Return complete output or explicit failure
\end{algorithmic}\end{algorithm}

The original CPU implementation uses a sparse bucket directory and per-wave proposal storage. Basic CUDA uses dense pending-directory scans, parallel frontier expansion, atomic proposal selection and launch barriers between expansion, commit and closure. It is a matched native implementation, not a claim of optimal GPU SSSP. Reset, scans, synchronization and output handoff have real costs. Finite iteration/work/capacity guards prevent an unbounded run; an aborted invocation is not a successful timing sample. Exhaustion or target completion after closure terminates a successful search. We assert neither native tie equivalence nor universal routing-quality preservation for this contract.
